\newtoggle{withbonus} \settoggle{withbonus}{false} % set true to show bonus points in grade table

% setting underlying course name (Grundlagenfach Mathematik, §-Unterricht, \ldots)
\newcommand{\mycourse}{Informatik (Nachprüfung)}
% name of the class
\newcommand{\myclass}{\faPeace\faIceCream~Klasse 2e~\faIceCream\faPeace}
% set date
\newcommand{\mydate}{Dienstag, 25. Juni 2024}

\title{\textbf{Theoretische Informatik: randomisierte Algorithmen}}
\author{Thomas Graf\\
\mycourse, \myclass}
\date{\mydate}
\maketitle
\thispagestyle{headandfoot}		

\begin{center}
	\fbox{\parbox{140mm}{\centering
			\begin{itemize}
				\item Dauer der Prüfung: $55$ Minuten
				%\item \textbf{Zeige in jeder Aufgabe unbedingt Deine Schritte!}
				\item erlaubte Hilfsmittel:
				\begin{itemize}
				    \item Schreibzeug
            \item \textbf{ohne Taschenrechner}
				\end{itemize}
	\end{itemize}}}
\end{center}
\vspace{10mm}

\begin{center}
	Vorname: \rule{45mm}{0.8pt}\qquad Nachname: \rule{45mm}{0.8pt}
\end{center}
\vspace{20mm}

\begin{center}
	\iftoggle{withbonus} {
		\combinedgradetable[h][questions]
	} {
		\gradetable[h][questions]
	}
\end{center}
\vspace{15mm}
\begin{center}
	Note: 
	\vspace{10mm}
	
	\rule{8mm}{1.2pt}
\end{center}
%\thispagestyle{empty}
\clearpage


\begin{questions}
% 2b
\question
\begin{parts}
  \part[1] Bestimmen Sie $\text{Nummer}(0101101)$.
  \part[1] Bestimmen Sie ein binäres Wort (ein binärer String) $x$, sodass $\text{Nummer}(x) = 82$.
  \part[1] Berechnen Sie $\abs{\text{Bin}(3099)}$. 
  \part
  Sei $y$ die binäre Darstellung ($2$-adische Darstellung) der drittgrössten Zahl, welche mit $9$ Bits dargestellt werden kann. 
  \begin{subparts}
    \subpart[1] Wie sieht $y$ aus?
    \subpart[1] Bestimmen Sie $\text{Nummer}\lr{y}$.
  \end{subparts}
  \part[1] Berechnen Sie (stellen Sie ohne Logarithmus dar) die Zahl $\ceil{\log_5\lr{89}}$.
  \part[1] Bestimmen Sie $\text{Prim(41)}$.
\end{parts}
\droptotalpoints


\question[3] Gegeben sind die beiden Datensätze
\begin{itemize}
  \item $x := 001001$ in Zürich und
  \item $y := 101010$ in Boston.
\end{itemize}
Berechnen Sie die exakte Fehlerwahrscheinlichkeit des Kommunikationsprotokolls in diesem Fall.


\question Gegeben seien zwei Datensätze $x$ und $y$ mit einer jeweiligen Länge von $10^{20}$ Bits. 
\begin{parts}
\part[2] Geben Sie eine vernünftige (möglichst kleine) obere Schranke für den Kommunikationsaufwand des Kommunikationsprotokolls zum Abgleich von $x$ und $y$ an.
\part[1] Wie viele Bits müsste ein deterministischer Algorithmus zur Beantwortung der Frage $x\stackrel{?}{=}y$ mindestens kommunizieren?
\end{parts}
\droptotalpoints


\question[2]
Für die Fehlerwahrscheinlichkeit $r$ des Protokolls haben wir die Abschätzung
\begin{align*}
    r \leq \frac{n-1}{\text{Prim}(n^2)}\leq \frac{n-1}{n^2/\ln(n^2)}\leq\frac{\ln(n^2)}{n}
\end{align*}
gefunden. Dabei ist $n$ die Länge der jeweiligen Datensätze (in Bits). Erklären Sie, warum wir im Protokoll eine Primzahl $p \leq n^2$ wählen und nicht lediglich eine Primzahl $p \leq n$.

\question[2]
Gegeben sind zwei Funktionen $f$ und $g$, welche auf ganz $\N$ definiert sind und $g(n) \neq 0$ für alle $n\in\N$. Es gelte die Gleichung
\begin{align*}
  \lim_{n\to\infty}\frac{f(n)}{g(n)} = 1.
\end{align*}
Ist es korrekt, dass daraus die Gleichung
\begin{align*}
  \lim_{n\to\infty} \abs{f(n) - g(n)} = 0
\end{align*}
folgt? Begründen Sie die Behauptung oder finden Sie ein Gegenbeispiel.
\end{questions}
